\section{A minimal two-learner/two-task model}

The following model is intentionally narrow. Learners $A,B$ and task regions $1,2$ have effective operational qualities
\begin{equation}
q_{iz}=c_{iz}-\lambda k_{iz}.
\label{eq:quality}
\end{equation}
Each learner has unit operational capacity, giving two divisions of labour:
\begin{equation}
D=(A\!\rightarrow\!1,B\!\rightarrow\!2),\qquad
X=(A\!\rightarrow\!2,B\!\rightarrow\!1).
\label{eq:divisions}
\end{equation}
Their values and difference are
\begin{equation}
R_D=q_{A1}+q_{B2},\quad R_X=q_{A2}+q_{B1},\quad
\Delta(C)=R_D-R_X.
\label{eq:division-values}
\end{equation}

\subsection{Assumptions and separability null case}

The model assumes two tasks, two learners, deterministic assignments, and no queues, abstention, mixed routing, transfer, interference, demand uncertainty, or shared training budget unless a later subsection explicitly introduces a coupling. It should not be read as a realistic HLS architecture.

Before adding capacity coupling, consider task-separable rewards and transitions:
\begin{equation}
R(C,a)=R_1(C_1,a_1)+R_2(C_2,a_2),\qquad
P(C'\mid C,a)=P_1(C'_1\mid C_1,a_1)P_2(C'_2\mid C_2,a_2).
\label{eq:separable}
\end{equation}
When no shared constraint or information couples the tasks, the Bellman maximisation factors. This null case is essential: a joint state notation alone does not create a joint competence-evolution problem.
